\begin{afterword}
\summaryhead

The post begins with a familiar cycle. A favoured theory, used to explain more and more, seems truer, so evidence against it is doubted and the theory is applied more widely. That cycle, the author says, is nothing compared with one that starts from a thought that feels good: a political system that will save the world, a great leader, a tonic that cures everything.

The mechanism is the halo effect, by which any perceived positive trait raises the perception of other positive traits. The author compares it to a fission chain reaction. Weak positive feeling is subcritical: each boost produces less than one further boost, and the effect dies out. Intense feeling attached to a ``Great Thingy'' reaches everywhere. A believing Communist sees the wisdom of Marx in every hamburger and every newspaper article; each interpretation confirms the idea and feels good, and the believer looks for more events to interpret. The author calls this an ``affective death spiral'' and promises a remedy in the next post.

\responsehead

The post names a process and describes it by analogy. It marks the analogy as one (``Metaphorically speaking''), and the physics it describes is stated correctly. What it does not supply is evidence that the process runs as described.

The step that turns the halo effect into a spiral is feedback: a judgment raised by the halo must raise the first judgment in turn, and the loop must grow. The post offers this step with ``perhaps'' and no evidence. The studies it relies on, summarized two days earlier in ``The Halo Effect,'' measure one trait coloring judgments of others on one occasion. None follows judgments feeding back over time. The chain reaction is the post's model of belief, and its key step is unsupported.

The central example does not show the mechanism the post introduces. The post says the familiar cycle of its first paragraph is ``nothing compared to'' the spiral that starts from good feeling, but it compares the two on no case. Its Communist, who finds Marx confirmed in every event, is mostly that first cycle: a theory used to interpret everything and confirmed by each use. What the example adds is pleasure, and the assertion that pleasure makes us want to believe. It contains no case of one positive trait raising the perception of another.

The example is also not new. Karl Popper described the same experience, down to the newspaper, in \textsc{Conjectures and Refutations}: a Marxist found confirming evidence ``on every page,'' and ``Whatever happened always confirmed it.'' Popper drew a lesson about testing: confirmations that come this easily are worth little. The post draws a lesson about feeling, and does not mention him.

In short: a named process, described by a chain-reaction analogy, whose feedback step is offered with ``perhaps'' and whose main example, close to one Popper gave, shows mostly the familiar confirmation cycle rather than the halo mechanism the post introduces.
\end{afterword}
